\paragraph{Purpose}
Policy gradient methods take noisy steps, and one step that is too large may ruin the
policy, after which the data collected is bad too. PPO (Schulman et al., 2017) lets us
\tB{reuse the same batch for several gradient steps while keeping the new policy close to the
policy that collected the data}, with a simple first-order objective.

\paragraph{Theory}
Let $\pi_{old}$ be the policy that generated the batch, and the probability ratio
\begin{center}
    $r_{t}(\bm{\theta})=\dfrac{\pi_{\bm{\theta}}(a_{t}|s_{t})}{\pi_{old}(a_{t}|s_{t})}$
\end{center}
(importance weight, $r_{t}=1$ at the start). The clipped objective to maximize is
\begin{center}
    \frB{$L^{CLIP}(\bm{\theta})=\mathbb{E}_{t}\left[\min\left(r_{t}(\bm{\theta})\hat{A}_{t},\
    \mathrm{clip}\left(r_{t}(\bm{\theta}),1-\varepsilon,1+\varepsilon\right)\hat{A}_{t}
    \right)\right]$}
\end{center}
with $\varepsilon\approx0.2$.
\begin{itemize}
    \item if $\hat{A}_{t}>0$ (good action), the objective grows with $r_{t}$ only up to
        $1+\varepsilon$: no incentive to make the action much more probable;
    \item if $\hat{A}_{t}<0$ (bad action), it is the symmetric case, the gradient vanishes
        once $r_{t}<1-\varepsilon$;
    \item the $\min$ makes the objective a \uB{pessimistic bound}: the clipping only removes
        the incentive to move further, it never hides a harmful move.
\end{itemize}

\paragraph{Algorithm}
\begin{enumerate}
    \item collect trajectories with $\pi_{old}$ and compute the advantages $\hat{A}_{t}$
        (critic and GAE);
    \item for several epochs, on mini-batches, maximize
        $L^{CLIP}-c_{1}\,L^{value}+c_{2}\,H\left[\pi_{\bm{\theta}}\right]$, with
        $L^{value}$ the squared error of the critic and $H$ an entropy bonus that keeps
        some exploration;
    \item set $\pi_{old}\leftarrow\pi_{\bm{\theta}}$ and repeat.
\end{enumerate}

\paragraph{Strengths and weaknesses}
\begin{itemize}
    \item simple, stable, and the default policy-gradient algorithm for a decade, including
        the training of language models from human feedback;
    \item it needs a critic (a second network of the size of the policy for LLMs) and
        several hyper-parameters, and it remains sensitive to implementation details.
\end{itemize}

\paragraph{From RL to language models} A language model is a policy: the \emph{state} is the
prompt and the tokens generated so far, the \emph{action} is the next token, the episode ends
at the end-of-text token, and a reward arrives for the complete answer
(see the alignment chapter).
